\documentclass{article}
\input{../../../lib.sty}

\title{\texttt{fn log\_alpha\_over\_alpha\_minus\_one\_lo}}
\author{Michael Shoemate}

\begin{document}
\maketitle

Proves soundness of \texttt{fn log\_alpha\_over\_alpha\_minus\_one\_lo} in \texttt{mod.rs}.

\section{Hoare Triple}
\subsection*{Precondition}
\texttt{alpha} $> 1$.

\subsection*{Pseudocode}
\lstinputlisting[language=Python,firstline=2,escapechar=|]{./pseudocode/log_alpha_over_alpha_minus_one_lo.py}

\subsection*{Postcondition}
\begin{theorem}
\label{postcondition}
For any setting of the input parameter \texttt{alpha} such that the given precondition holds,
\texttt{log\_alpha\_over\_alpha\_minus\_one\_lo} either returns \texttt{Err(e)},
or returns \texttt{Ok(out)} where
\[
\texttt{out} \le \log\!\left(\frac{\texttt{alpha}}{\texttt{alpha} - 1}\right).
\]
\end{theorem}

\section{Proof}
Assume the precondition is met.

\begin{lemma}
\label{err-e}
\texttt{log\_alpha\_over\_alpha\_minus\_one\_lo} only returns \texttt{Err(e)} if one of the
fallible helper calls in the pseudocode returns \texttt{Err(e)}.
\end{lemma}

\begin{proof}
The only fallible operations are the calls on lines \ref{line:div_lo} and \ref{line:ln1p_hi}.
\end{proof}

\begin{lemma}
\label{ok-out}
If the outcome of \texttt{log\_alpha\_over\_alpha\_minus\_one\_lo} is \texttt{Ok(out)}, then
\[
\texttt{out} \le \log\!\left(\frac{\texttt{alpha}}{\texttt{alpha} - 1}\right).
\]
\end{lemma}

\begin{proof}
By the proof definition of \texttt{neg\_inf\_div},
\[
\texttt{inv\_alpha\_lo} \le \frac{1}{\texttt{alpha}}.
\]
Multiplying by $-1$ reverses the inequality:
\[
\texttt{neg\_inv\_alpha\_hi} = -\texttt{inv\_alpha\_lo} \ge -\frac{1}{\texttt{alpha}}.
\]
Since \texttt{alpha} $> 1$, both arguments lie in $(-1,0)$, so by the proof definition of
\texttt{inf\_ln\_1p},
\[
\texttt{ln\_one\_minus\_inv\_alpha\_hi}
\ge \ln\!\left(1 - \frac{1}{\texttt{alpha}}\right).
\]
Multiplying by $-1$ gives
\[
\texttt{out}
=
-\texttt{ln\_one\_minus\_inv\_alpha\_hi}
\le
-\ln\!\left(1 - \frac{1}{\texttt{alpha}}\right)
= \log\!\left(\frac{\texttt{alpha}}{\texttt{alpha} - 1}\right).
\]
\end{proof}

\begin{proof}[Proof of Theorem~\ref{postcondition}]
Holds by Lemma~\ref{err-e} and Lemma~\ref{ok-out}.
\end{proof}

\end{document}
